\documentclass[../main.tex]{subfiles}
\begin{document}

1. \textbf{Kết quả đạt được và Kiến thức tích lũy}
Qua quá trình thực hiện đồ án thiết kế và xây dựng hệ thống quản lý thực phẩm thông minh NaviMart, nhóm chúng em đã gặt hái được nhiều kết quả quan trọng và tích lũy thêm nhiều kinh nghiệm quý báu:

\begin{itemize}
    \item \textbf{Về quy trình phát triển phần mềm:} Nhóm đã hiểu rõ và áp dụng thành công quy trình xây dựng phần mềm từ công đoạn khảo sát, phân tích yêu cầu cho đến thiết kế giao diện, xây dựng kiến trúc, kiểm thử hệ thống và triển khai thực tế theo tiêu chuẩn kỹ năng nghề nghiệp ITSS.
    \item \textbf{Về kỹ năng kỹ thuật:} Các thành viên được nâng cao năng lực lập trình chuyên sâu, đồng thời làm chủ các công nghệ hiện đại tiêu biểu như NestJS cho hệ thống máy chủ và React cho giao diện người dùng. Nhóm cũng hiểu sâu sắc cách thiết kế và giao tiếp dữ liệu thông qua chuẩn API cũng như các kết nối thời gian thực.
    \item \textbf{Về kỹ năng mềm và quản lý dự án:} Kỹ năng làm việc nhóm, phân chia công việc và trau dồi kỹ năng biểu diễn các mô hình UML trực quan được rèn luyện xuyên suốt quá trình thực hiện sản phẩm.
\end{itemize}

2. \textbf{Hạn chế và Hướng phát triển tương lai}
Mặc dù hệ thống NaviMart đã hoàn thiện các chức năng cốt lõi giúp người dùng quản lý kho thực phẩm và lên thực đơn dễ dàng, ứng dụng do nhóm thiết kế và xây dựng vẫn còn một số điểm cần tiếp tục hoàn thiện. 

Trong thời gian tới, nhóm dự định sẽ tiếp tục nâng cấp sản phẩm với các định hướng sau:
\begin{itemize}
    \item \textbf{Tối ưu hóa trợ lý thông minh:} Nâng cấp khả năng phân tích dữ liệu để đưa ra những gợi ý công thức nấu ăn mang tính cá nhân hóa cao hơn, dựa trên thói quen dinh dưỡng cũng như tình trạng sức khỏe của từng thành viên trong gia đình.
    \item \textbf{Đa nền tảng và thiết bị:} Phát triển thêm phiên bản ứng dụng di động gốc chạy trên hệ điều hành điện thoại nhằm tận dụng tối đa phần cứng máy ảnh, giúp quá trình quét mã vạch nhận diện thực phẩm diễn ra trơn tru và chính xác nhất.
    \item \textbf{Mở rộng tính năng cộng đồng:} Xây dựng nền tảng cho phép người dùng chia sẻ các công thức nấu ăn tự sáng tạo lên hệ thống chung, từ đó phát triển một cộng đồng chia sẻ tri thức ẩm thực bền vững và gắn kết hơn.
\end{itemize}

\end{document}